
\documentclass[11pt,a4paper]{article}
\usepackage[margin=1in]{geometry}
\usepackage{amsmath,amssymb}
\usepackage{graphicx}
\usepackage{booktabs}
\usepackage{longtable}
\usepackage{hyperref}
\usepackage{placeins}
\usepackage{xcolor}
\usepackage{caption}
\usepackage{subcaption}
\usepackage{enumitem}
\hypersetup{colorlinks=true,linkcolor=blue,urlcolor=blue,citecolor=blue}
\newcommand{\code}[1]{\texttt{#1}}
\newcommand{\maybeincludegraphics}[2][]{%
\IfFileExists{#2}{\includegraphics[#1]{#2}}{\fbox{\parbox[c][0.34\textheight][c]{0.88\linewidth}{\centering Missing figure: \code{#2}\\Run the MATLAB script in this report directory.}}}%
}

\title{High-Order Explicit Manufactured FDA Solver\\Verification Report}
\author{Generated report for \code{highOrderManufacturedFDAExplicit}}
\date{\today}
\begin{document}
\maketitle

\begin{abstract}
This report documents the manufactured-solution finite-volume solver \code{highOrderManufacturedFDAExplicit}. The solver advances a monodomain-like activation equation coupled to the local state variables $u_1$, $u_2$, and $u_3$. The spatial discretization may use standard OpenFOAM finite volumes or high-order local reconstruction by LRE. Time integration is explicit forward Euler. The convergence figures are generated by \code{plot\_convergence.m} from the data in \code{/home/pablo/OpenFOAM/pablo-v2312/run/tutorials\_electro/highOrderManufacturedFDAExplicit/results/example0/2D}.
\end{abstract}

\section{Manufactured FDA problem}
The solved variables are the transmembrane potential $V_m$ and the state variables $u_1$, $u_2$, and $u_3$. In the manufactured problem,
\begin{align}
\frac{\partial V_m}{\partial t} &= \frac{1}{\chi C_m}\nabla\cdot(D\nabla V_m) - I_{\mathrm{ion}}(V_m,u_1,u_2,u_3),\\
\frac{\partial u_1}{\partial t} &= (u_1+u_3-V_m)^2u_2^2 + \frac{1}{2}(u_1+u_3-V_m)u_2^2(V_m-u_3),\\
\frac{\partial u_2}{\partial t} &= -(u_1+u_3-V_m)u_2^3,\\
\frac{\partial u_3}{\partial t} &= 0.
\end{align}
The ionic current used in the PDE is
\begin{equation}
I_{\mathrm{ion}} = -\frac{1}{2}(u_1+u_3-V_m)u_2^2(V_m-u_3)
+ \frac{\beta}{\chi C_m}(V_m-u_3),
\end{equation}
where $\beta$ is the exact Laplacian eigenvalue induced by the conductivity tensor and by the selected manufactured spatial mode.

For the two-dimensional runs reported here,
\begin{align}
F(x,y) &= \cos(\pi x)\cos(2\pi y),\\
G(x,y) &= 1 + x y^2,\\
V_m^{\star}(x,y,t) &= \sqrt{1+t}\,F(x,y),\\
u_1^{\star}(x,y,t) &= (1+t)G(x,y) + V_m^{\star}(x,y,t),\\
u_2^{\star}(x,y,t) &= \frac{1}{(1+t)\sqrt{G(x,y)}},\\
u_3^{\star} &= 0.
\end{align}
The code also contains analogous one- and three-dimensional definitions.

\section{Spatial discretization}
The potential equation can use either the standard OpenFOAM finite-volume Laplacian or a high-order flux constructed from LRE. For the high-order option, local Taylor reconstructions are built from cell stencils. The reconstruction uses first derivatives for \code{p1}, first and second derivatives for \code{p2}, and up to third derivatives for \code{p3}. Face quadrature is used to integrate the diffusive flux, and cell quadrature is used to volume-average the nonlinear ionic current when high-order state reconstruction is enabled.

The switches \code{useHighOrder\_Vm} and \code{useHighOrder\_states} allow the potential and state variables to use different high-order levels. The expected spatial orders used in the plots are $2$ for \code{p1}, $3$ for \code{p2}, and $4$ for \code{p3}. The standard non-HO option is shown with the second-order guide line.

\section{Explicit time integration}
At each time step the solver computes the ionic current and the Laplacian from known values at $t^n$, then applies forward Euler,
\begin{align}
V_m^{n+1} &= V_m^n + \Delta t\left[\frac{1}{\chi C_m}L_h(V_m^n)-I_{\mathrm{ion},h}(V_m^n,u^n)\right],\\
u_i^{n+1} &= u_i^n + \Delta t\,R_i(V_m^n,u^n),\qquad i=1,2,3.
\end{align}
The manufactured exact values are imposed at the relevant Dirichlet boundaries after each update.

\section{Error measures and convergence rates}
The solver writes cell-centred volume-weighted $L_1$, $L_2$, and $L_\infty$ errors for $V_m$, $u_1$, and $u_2$. It also writes Gauss-reconstructed errors, denoted here by the superscript $G$, obtained by evaluating the reconstructed field at quadrature points. Pairwise rates use consecutive meshes,
\begin{equation}
p_{N_1,N_2}=\frac{\log(E_{N_1}/E_{N_2})}{\log(h_{N_1}/h_{N_2})}, \qquad h_N = 1/N.
\end{equation}
The summary table reports the all-mesh log--log fit of the cell-centred $L_2$ errors.

\section{Observed orders}
\begin{longtable}{p{0.38\linewidth}lrrrrrr}
\caption{Explicit solver: observed cell-centred $L_2$ convergence orders obtained from the all-mesh log--log fit. The expected orders use $p1\mapsto2$, $p2\mapsto3$, $p3\mapsto4$; the non-HO standard option is treated as second order.}\label{tab:explicit-orders}\\
\toprule
Case & $\Delta t$ & Exp. $V_m$ & $V_m$ & Exp. states & $u_1$ & $u_2$ & $V_m^G$ \\
\midrule
\endfirsthead
\toprule
Case & $\Delta t$ & Exp. $V_m$ & $V_m$ & Exp. states & $u_1$ & $u_2$ & $V_m^G$ \\
\midrule
\endhead
hoVm\_NO\_hoStates\_NO & $10^{-3}$ & 2 & 1.834 & 2 & 0.389 & 0.007 & 2.714 \\
hoVm\_NO\_hoStates\_NO & $10^{-4}$ & 2 & 1.979 & 2 & 1.885 & 0.264 & 2.859 \\
hoVm\_NO\_hoStates\_NO & $10^{-5}$ & 2 & 1.995 & 2 & 2.037 & 1.336 & 2.176 \\
hoVm\_p1\_hoStates\_NO & $10^{-3}$ & 2 & 1.992 & 2 & 1.586 & 0.255 & 2.145 \\
hoVm\_p1\_hoStates\_NO & $10^{-4}$ & 2 & 1.994 & 2 & 1.703 & 1.355 & 2.148 \\
hoVm\_p1\_hoStates\_NO & $10^{-5}$ & 2 & 1.994 & 2 & 1.703 & 1.705 & 2.149 \\
hoVm\_p1\_hoStates\_p1 & $10^{-3}$ & 2 & 1.992 & 2 & 1.586 & 0.255 & 2.145 \\
hoVm\_p1\_hoStates\_p1 & $10^{-4}$ & 2 & 1.994 & 2 & 1.703 & 1.355 & 2.148 \\
hoVm\_p1\_hoStates\_p1 & $10^{-5}$ & 2 & 1.994 & 2 & 1.703 & 1.705 & 2.149 \\
hoVm\_p2\_hoStates\_NO & $10^{-3}$ & 3 & 2.069 & 2 & 1.317 & 0.098 & 2.514 \\
hoVm\_p2\_hoStates\_NO & $10^{-4}$ & 3 & 2.118 & 2 & 2.229 & 0.877 & 2.557 \\
hoVm\_p2\_hoStates\_NO & $10^{-5}$ & 3 & 2.123 & 2 & 2.151 & 1.969 & 2.562 \\
hoVm\_p2\_hoStates\_p1 & $10^{-3}$ & 3 & 2.139 & 2 & 1.172 & 0.071 & 2.658 \\
hoVm\_p2\_hoStates\_p1 & $10^{-4}$ & 3 & 2.210 & 2 & 2.296 & 0.766 & 2.711 \\
hoVm\_p2\_hoStates\_p1 & $10^{-5}$ & 3 & 2.217 & 2 & 2.245 & 1.973 & 2.717 \\
hoVm\_p2\_hoStates\_p2 & $10^{-3}$ & 3 & 2.140 & 3 & 1.156 & 0.078 & 2.658 \\
hoVm\_p2\_hoStates\_p2 & $10^{-4}$ & 3 & 2.211 & 3 & 2.336 & 0.773 & 2.712 \\
hoVm\_p2\_hoStates\_p2 & $10^{-5}$ & 3 & 2.219 & 3 & 2.263 & 1.953 & 2.718 \\
hoVm\_p3\_hoStates\_NO & $10^{-3}$ & 4 & 1.953 & 2 & 0.437 & 0.009 & 2.382 \\
hoVm\_p3\_hoStates\_NO & $10^{-4}$ & 4 & 2.143 & 2 & 1.909 & 0.282 & 2.573 \\
hoVm\_p3\_hoStates\_NO & $10^{-5}$ & 4 & 2.166 & 2 & 2.235 & 1.349 & 2.596 \\
hoVm\_p3\_hoStates\_p1 & $10^{-3}$ & 4 & 2.579 & 2 & 0.446 & -0.007 & 3.168 \\
hoVm\_p3\_hoStates\_p1 & $10^{-4}$ & 4 & 2.297 & 2 & 1.471 & 0.139 & 2.900 \\
hoVm\_p3\_hoStates\_p1 & $10^{-5}$ & 4 & 2.263 & 2 & 2.214 & 1.164 & 2.865 \\
hoVm\_p3\_hoStates\_p3 & $10^{-3}$ & 4 & 3.281 & 4 & 0.438 & 0.004 & 3.808 \\
hoVm\_p3\_hoStates\_p3 & $10^{-4}$ & 4 & 2.393 & 4 & 1.506 & 0.181 & 3.037 \\
hoVm\_p3\_hoStates\_p3 & $10^{-5}$ & 4 & 2.341 & 4 & 2.245 & 1.129 & 2.983 \\
\bottomrule
\end{longtable}


\section{Convergence figures}
Run \code{plot\_convergence.m} from this report directory to create the PDF and PNG figures. The plots compare all cases available in the explicit two-dimensional result directory for $\Delta t=10^{-3}$, $10^{-4}$, and $10^{-5}$.

\begin{figure}[htbp]
\centering
\maybeincludegraphics[width=0.92\linewidth]{figures/explicit_Vm_dt_1p0em03.pdf}
%\caption{Explicit solver, Vm cell-centred $L_2$ error for $\Delta t=$10^{-3}$$. Dotted guide lines show $O(h^2)$, $O(h^3)$, and $O(h^4)$.}
\label{fig:explicit-Vm-1p0em03}
\end{figure}

\begin{figure}[htbp]
\centering
\maybeincludegraphics[width=0.92\linewidth]{figures/explicit_u1_dt_1p0em03.pdf}
%\caption{Explicit solver, u1 cell-centred $L_2$ error for $\Delta t=$10^{-3}$$. Dotted guide lines show $O(h^2)$, $O(h^3)$, and $O(h^4)$.}
\label{fig:explicit-u1-1p0em03}
\end{figure}

\begin{figure}[htbp]
\centering
\maybeincludegraphics[width=0.92\linewidth]{figures/explicit_u2_dt_1p0em03.pdf}
%\caption{Explicit solver, u2 cell-centred $L_2$ error for $\Delta t=$10^{-3}$$. Dotted guide lines show $O(h^2)$, $O(h^3)$, and $O(h^4)$.}
\label{fig:explicit-u2-1p0em03}
\end{figure}

\begin{figure}[htbp]
\centering
\maybeincludegraphics[width=0.92\linewidth]{figures/explicit_Vm_dt_1p0em04.pdf}
%\caption{Explicit solver, Vm cell-centred $L_2$ error for $\Delta t=$10^{-4}$$. Dotted guide lines show $O(h^2)$, $O(h^3)$, and $O(h^4)$.}
\label{fig:explicit-Vm-1p0em04}
\end{figure}

\begin{figure}[htbp]
\centering
\maybeincludegraphics[width=0.92\linewidth]{figures/explicit_u1_dt_1p0em04.pdf}
%\caption{Explicit solver, u1 cell-centred $L_2$ error for $\Delta t=$10^{-4}$$. Dotted guide lines show $O(h^2)$, $O(h^3)$, and $O(h^4)$.}
\label{fig:explicit-u1-1p0em04}
\end{figure}

\begin{figure}[htbp]
\centering
\maybeincludegraphics[width=0.92\linewidth]{figures/explicit_u2_dt_1p0em04.pdf}
%\caption{Explicit solver, u2 cell-centred $L_2$ error for $\Delta t=$10^{-4}$$. Dotted guide lines show $O(h^2)$, $O(h^3)$, and $O(h^4)$.}
\label{fig:explicit-u2-1p0em04}
\end{figure}

\begin{figure}[htbp]
\centering
\maybeincludegraphics[width=0.92\linewidth]{figures/explicit_Vm_dt_1p0em05.pdf}
%\caption{Explicit solver, Vm cell-centred $L_2$ error for $\Delta t=$10^{-5}$$. Dotted guide lines show $O(h^2)$, $O(h^3)$, and $O(h^4)$.}
\label{fig:explicit-Vm-1p0em05}
\end{figure}

\begin{figure}[htbp]
\centering
\maybeincludegraphics[width=0.92\linewidth]{figures/explicit_u1_dt_1p0em05.pdf}
%\caption{Explicit solver, u1 cell-centred $L_2$ error for $\Delta t=$10^{-5}$$. Dotted guide lines show $O(h^2)$, $O(h^3)$, and $O(h^4)$.}
\label{fig:explicit-u1-1p0em05}
\end{figure}

\begin{figure}[htbp]
\centering
\maybeincludegraphics[width=0.92\linewidth]{figures/explicit_u2_dt_1p0em05.pdf}
%\caption{Explicit solver, u2 cell-centred $L_2$ error for $\Delta t=$10^{-5}$$. Dotted guide lines show $O(h^2)$, $O(h^3)$, and $O(h^4)$.}
\label{fig:explicit-u2-1p0em05}
\end{figure}


\section{Notes}
The explicit results may be affected by temporal error unless $\Delta t$ is sufficiently small for the spatial error to dominate. When interpreting the plots, compare the three time-step levels and the reconstructed $V_m^G$ column in Table~\ref{tab:explicit-orders} against the cell-centred $V_m$ order.

\end{document}
